\chapter{Markov Chain Steady States and Mixing}

In the previous chapter you met the Markov chain: a system that hops between
states, where the next step depends only on where you are now. We classified
states as recurrent or transient and computed multi-step matrix powers $P^n$.
But what happens as time heads toward infinity?

If you let a Markov chain run for a long time, it very often \emph{forgets where
it started}. The fraction of time it spends in each state settles down to a
fixed set of proportions, regardless of the initial state. Individual walkers
keep hopping, yet the overall distribution becomes motionless.

This chapter is about that settling-down. We ask three questions and answer
all of them:

\begin{itemize}
  \item \textbf{Is there a long-run distribution at all?} (Stationary
    distributions, and how to find them as a left eigenvector.)
  \item \textbf{Does the chain actually converge to it, no matter where it
    starts?} (The convergence theorem for ergodic chains.)
  \item \textbf{How long do you have to wait?} (Mixing time.)
\end{itemize}

Along the way we meet \emph{reversibility} and detailed balance --- an algebraic condition
that identifies the steady state directly --- and examine two major applications:
Google's PageRank, and Markov chain Monte Carlo (MCMC), the computational engine of
modern Bayesian statistics.

\chapterline

\section{The Long Run: A First Look}
\label{ch24:longrun}

Let us anchor everything in a tiny, concrete world. Imagine the weather in a town
that has only two moods: \textbf{Sunny} (state $S$) and \textbf{Rainy} (state
$R$). The weather tomorrow depends only on the weather today. Suppose:

\begin{itemize}
  \item If it is sunny today, tomorrow is sunny with probability $0.8$ and rainy
    with probability $0.2$.
  \item If it is rainy today, tomorrow is rainy with probability $0.6$ and sunny
    with probability $0.4$.
\end{itemize}

We collect these into a \textbf{transition matrix} $P$, where the entry
$P_{ij}$ is the probability of moving from state $i$ to state $j$:
\[
  P = \begin{pmatrix} P_{SS} & P_{SR} \\ P_{RS} & P_{RR} \end{pmatrix}
    = \begin{pmatrix} 0.8 & 0.2 \\ 0.4 & 0.6 \end{pmatrix}.
\]
Each row sums to $1$, because from any state \emph{something} must happen
tomorrow. (A matrix with nonnegative entries and rows summing to one is called a
\textbf{stochastic matrix}; every Markov chain is one of these and vice versa.)

\begin{figure}[H]
  \centering
  \begin{tikzpicture}[>=Stealth, node distance=3.4cm]
    \node[circle,draw=sourceblue,very thick,fill=sourceblue!12,minimum size=1.2cm] (S) {$S$};
    \node[circle,draw=targetgreen,very thick,fill=targetgreen!12,minimum size=1.2cm,right=of S] (R) {$R$};
    \draw[->,thick,bend left=18] (S) to node[above]{$0.2$} (R);
    \draw[->,thick,bend left=18] (R) to node[below]{$0.4$} (S);
    \draw[->,thick] (S) edge[loop left,looseness=6] node[left]{$0.8$} (S);
    \draw[->,thick] (R) edge[loop right,looseness=6] node[right]{$0.6$} (R);
  \end{tikzpicture}
  \caption{The two-state weather chain. Arrows are transition probabilities;
  the loops are the chances of staying put.}
  \label{ch24:fig-weather-diagram}
\end{figure}

Now matrix multiplication reveals its predictive power. Suppose today the
weather is described by a \textbf{row vector} of probabilities,
$\pi^{(0)} = (\,\P(\text{sunny}),\ \P(\text{rainy})\,)$. What is the distribution
tomorrow? You multiply on the right by $P$:
\[
  \pi^{(1)} = \pi^{(0)} P.
\]
Why on the right, and why a row vector? Because the probability of being sunny
tomorrow is ``probability sunny today times $P_{SS}$, plus probability rainy
today times $P_{RS}$'' --- exactly the dot product of the row vector with the
$S$-column of $P$. Iterating, after $n$ days,
\[
  \pi^{(n)} = \pi^{(0)} P^n.
\]
The whole dynamics of the chain is hidden inside the powers of one matrix.

\begin{walkbox}
Start certain that today is sunny, so $\pi^{(0)} = (1,0)$. Compute the
distribution for the next few days. Watch what happens.
\[
  \begin{aligned}
    \pi^{(1)} &= (1,0)P = (0.8,\,0.2),\\
    \pi^{(2)} &= (0.8,0.2)P = (0.8\cdot0.8 + 0.2\cdot0.4,\ \ 0.8\cdot0.2 + 0.2\cdot0.6) = (0.72,\,0.28),\\
    \pi^{(3)} &= (0.72,0.28)P = (0.688,\,0.312),\\
    \pi^{(4)} &\approx (0.6752,\,0.3248),\qquad
    \pi^{(5)} \approx (0.6701,\,0.3299).
  \end{aligned}
\]
Now start \emph{certain that today is rainy}, $\pi^{(0)}=(0,1)$:
\[
  \pi^{(1)} = (0.4,0.6),\ \ \pi^{(2)}=(0.56,0.44),\ \ \pi^{(3)}=(0.624,0.376),\ \ \pi^{(5)}\approx(0.6598,0.3402).
\]
The two sequences are marching toward the \emph{same} pair of numbers,
$(2/3,\,1/3)$, from opposite sides. The chain is forgetting where it began.
\end{walkbox}

Figure~\ref{ch24:fig-weather-converge} plots both runs. This convergence to a
single distribution, regardless of the start, is the central phenomenon of the
chapter. Our job now is to (a) name the limit, (b) compute it without grinding
out powers, and (c) understand exactly when it is guaranteed.

\begin{figure}[H]
  \centering
  \begin{tikzpicture}
    \begin{axis}[
        width=12cm, height=6.2cm,
        xlabel={day $n$}, ylabel={$\P(\text{sunny})$ on day $n$},
        xmin=0, xmax=10, ymin=0, ymax=1.05,
        title style={font=\small},
        title={Convergence of marginal distribution to stationary state $\pi = (2/3, 1/3)$},
        grid=both, grid style={gray!18},
        legend pos=south east, legend cell align=left,
      ]
      \addplot[sourceblue,very thick,mark=*,mark size=1.4pt] coordinates {
        (0,1)(1,0.8)(2,0.72)(3,0.688)(4,0.6752)(5,0.67008)(6,0.668032)(7,0.6672128)(8,0.66688512)(9,0.666754048)(10,0.6667016192)};
      \addlegendentry{start sunny}
      \addplot[densityred,very thick,mark=square*,mark size=1.4pt] coordinates {
        (0,0)(1,0.4)(2,0.56)(3,0.624)(4,0.6496)(5,0.65984)(6,0.663936)(7,0.6655744)(8,0.66622976)(9,0.666491904)(10,0.6665967616)};
      \addlegendentry{start rainy}
      \addplot[black,dashed,thick,domain=0:10] {0.66666667};
      \addlegendentry{$\pi_S = 2/3$}
    \end{axis}
  \end{tikzpicture}
  \caption{Both starting conditions converge to the same stationary
  probability of sun, $2/3$, within a handful of days. The gap between the curves
  --- a memory of the initial state --- shrinks geometrically.}
  \label{ch24:fig-weather-converge}
\end{figure}
\section{Stationary Distributions}
\label{ch24:stationary}
The limit we just watched has a defining property: once you are \emph{in} it,
applying $P$ leaves you unchanged. It is a fixed point of the dynamics.
\begin{definition}[Stationary distribution]
\label{ch24:def-stationary}
A probability row vector $\pi = (\pi_1,\dots,\pi_m)$ (so $\pi_i \ge 0$ and
$\sum_i \pi_i = 1$) is a \textbf{stationary distribution} (also called an
\emph{invariant} or \emph{equilibrium} distribution) for a chain with transition
matrix $P$ if
\[
  \pi P = \pi.
\]
Equivalently, for every state $j$, \quad $\displaystyle \pi_j = \sum_i \pi_i\,P_{ij}$.
\end{definition}
Read the componentwise version as a \emph{conservation law}. Picture a large
population of independent walkers spread over the states in proportions $\pi$.
In one step, each state $j$ \emph{receives} a flow $\pi_i P_{ij}$ from every state
$i$. The distribution is stationary exactly when the total inflow into $j$ equals
the amount $\pi_j$ that was there to begin with: probability mass sloshes around,
but the overall proportions never change.
\begin{intuitionbox}
``Stationary'' does not mean nothing moves. Individual walkers continue to transition
between states at each step. What remains constant is the aggregate proportion of probability
mass residing in each state. Equilibrium is dynamic rather than static.
\end{intuitionbox}
\subsection{Finding it by hand}
For the weather chain, $\pi P = \pi$ with $\pi = (\pi_S, \pi_R)$ gives
\[
  \begin{cases}
    0.8\,\pi_S + 0.4\,\pi_R = \pi_S,\\
    0.2\,\pi_S + 0.6\,\pi_R = \pi_R.
  \end{cases}
\]
Both equations reduce to the same relation $0.2\,\pi_S = 0.4\,\pi_R$, i.e.
$\pi_S = 2\pi_R$. Because the rows of $P$ sum to one, the
columns of $P - I$ sum to zero, which means one of the balance equations is always
linearly dependent. The missing constraint is supplied by the probability normalization
$\pi_S + \pi_R = 1$. Solving,
\[
  2\pi_R + \pi_R = 1 \ \Rightarrow\ \pi_R = \tfrac13,\qquad \pi_S = \tfrac23.
\]
So $\pi = (2/3,\,1/3)$ --- precisely the numbers Figure~\ref{ch24:fig-weather-converge}
crept toward. In the long run this town is sunny two days out of three.
\begin{trapbox}
A stationary distribution must be \emph{normalized}. The equation $\pi P = \pi$ is
homogeneous: if $\pi$ solves it, so does $5\pi$ and so does $0$. Those are not
distributions. You always need the extra condition $\sum_i \pi_i = 1$ to pin down
the actual probability vector. In practice: solve $\pi P = \pi$ for the
\emph{ratios} between components, then scale so they add to $1$.
\end{trapbox}

\begin{recoverybox}
\textbf{Pause and think: doubly stochastic matrices.} What happens if every column of $P$ also sums to $1$ (a \emph{doubly stochastic} matrix, where total inflow capacity equals total outflow capacity everywhere)?
\emph{Ten-second thought:} Try the uniform distribution $\pi = (1/m, \dots, 1/m)$. Then the $j$-th entry of $\pi P$ is $\sum_i \frac{1}{m} P_{ij} = \frac{1}{m}\sum_i P_{ij} = \frac{1}{m}\cdot 1 = \pi_j$. So for any doubly stochastic chain, the uniform distribution is automatically stationary: every state receives equal long-run weight.
\end{recoverybox}

\subsection{It really is a left eigenvector}
\begin{connectionbox}
The equation $\pi P = \pi$ is the eigenvalue equation $\pi P = 1\cdot\pi$. A standard (right) eigenvector $v$ satisfies $P v = \lambda v$: the matrix acts on a column vector by stretching it without rotating its direction. A \textbf{left eigenvector} is a row vector $\pi$ satisfying $\pi P = \lambda \pi$: when you mix the rows of $P$ using weights $\pi$, the resulting mixture is just a scaled copy of $\pi$ itself. \emph{A stationary distribution is exactly a left eigenvector of $P$ for the eigenvalue $\lambda = 1$, normalized to sum to one.}

Why is $1$ always an eigenvalue? Because $P$ has constant row sums equal to $1$,
the all-ones \emph{column} vector $\mathbf{1}$ satisfies $P\mathbf{1} = \mathbf{1}$
--- it is a \emph{right} eigenvector with eigenvalue $1$. A matrix and its
transpose have the same eigenvalues, so $1$ is a left eigenvalue too, guaranteeing
a (possibly complex, but in fact real and nonnegative) left eigenvector $\pi$.
That existence is half of the Perron--Frobenius theorem --- the foundational result in
linear algebra guaranteeing that any matrix with nonnegative entries and connected paths
has a dominant eigenvalue equal to $1$ whose eigenvector can be chosen strictly positive.
The companion linear algebra book treats it in full. The key fact we borrow: for a
``nice'' (irreducible, aperiodic) stochastic matrix, the eigenvalue $1$ is
\emph{simple} (multiplicity one), every other eigenvalue has modulus strictly
less than one, and the eigenvector $\pi$ has all entries strictly positive.
\end{connectionbox}
The eigenvalue spectrum controls the \emph{rate} of
convergence, as we will see in Section~\ref{ch24:mixing}. The second-largest
eigenvalue is the rate at which memory of the starting state decays.
\begin{example}[A three-state chain]
\label{ch24:ex-three-state}
Consider three web pages $A$, $B$, $C$ with a surfer who clicks links:
\[
  P = \begin{pmatrix}
    0 & 1 & 0 \\
    \tfrac14 & 0 & \tfrac34 \\
    \tfrac12 & \tfrac12 & 0
  \end{pmatrix}.
\]
\begin{figure}[H]
  \centering
  \begin{tikzpicture}[>=Stealth, node distance=3.2cm]
    \node[circle,draw=sourceblue,very thick,fill=sourceblue!12,minimum size=1.1cm] (A) {$A$};
    \node[circle,draw=targetgreen,very thick,fill=targetgreen!12,minimum size=1.1cm,right=of A] (B) {$B$};
    \node[circle,draw=quotientorange,very thick,fill=quotientorange!16,minimum size=1.1cm,below=2.4cm of $(A)!0.5!(B)$] (C) {$C$};
    \draw[->,thick,bend left=15] (A) to node[above]{$1$} (B);
    \draw[->,thick,bend left=15] (B) to node[below,pos=0.4]{$\tfrac14$} (A);
    \draw[->,thick,bend left=15] (B) to node[right]{$\tfrac34$} (C);
    \draw[->,thick,bend left=15] (C) to node[left]{$\tfrac12$} (A);
    \draw[->,thick,bend left=15] (C) to node[pos=0.45,left=2pt]{$\tfrac12$} (B);
  \end{tikzpicture}
  \caption{The three-page surfer chain of Example~\ref{ch24:ex-three-state}.
  Page $A$ sends all its weight to $B$; $B$ splits between $A$ and $C$; $C$ splits
  between $A$ and $B$. Page $B$ receives all of $A$'s outbound weight, which is why it
  wins the largest stationary share.}
  \label{ch24:fig-three-state-diagram}
\end{figure}
We seek $\pi = (a,b,c)$ with $\pi P = \pi$:
\[
  \begin{aligned}
    a &= \tfrac14 b + \tfrac12 c, \\
    b &= a + \tfrac12 c, \\
    c &= \tfrac34 b.
  \end{aligned}
\]
Substitute $c = \tfrac34 b$ into the first: $a = \tfrac14 b + \tfrac12\cdot\tfrac34 b
= \tfrac14 b + \tfrac38 b = \tfrac58 b$. Check the second:
$a + \tfrac12 c = \tfrac58 b + \tfrac38 b = b$, consistent. Normalizing,
\[
  a + b + c = \tfrac58 b + b + \tfrac34 b = \tfrac{19}{8} b = 1
  \ \Rightarrow\ b = \tfrac{8}{19}.
\]
Hence $\pi = \bigl(\tfrac{5}{19},\,\tfrac{8}{19},\,\tfrac{6}{19}\bigr)$. Page $B$
is most visited: as Figure~\ref{ch24:fig-three-state-diagram} shows, it is the
page receiving the largest net incoming weight ($1$ from $A$ and $\tfrac12$ from $C$).
\end{example}
\section{Does the Chain Converge? The Ergodic Theorem}
\label{ch24:ergodic}
Existence of a stationary $\pi$ does not by itself guarantee that the chain
\emph{approaches} it. Two failure modes block convergence, and naming them tells
us exactly what to require.
\subsection{Failure mode 1: the chain is reducible}
If the state space splits into pieces that cannot reach one another, the long-run
behavior depends on which piece you started in --- there is no single universal
limit. We need every state to be reachable from every other.
\begin{definition}[Irreducible]
\label{ch24:def-irreducible}
A chain is \textbf{irreducible} if for every pair of states $i,j$ there is some
number of steps $n$ with $(P^n)_{ij} > 0$: you can get from anywhere to anywhere,
eventually. (On the diagram: the directed graph of allowed transitions is
strongly connected.)
\end{definition}
\subsection{Failure mode 2: the chain is periodic}
Even an irreducible chain can fail to settle if it marches in lockstep. Consider
two states with $P = \left(\begin{smallmatrix} 0 & 1 \\ 1 & 0\end{smallmatrix}\right)$:
it flips $A \to B \to A \to B \to \cdots$ forever. The stationary distribution is
$(1/2,1/2)$, but starting at $A$ the distribution is $(1,0),(0,1),(1,0),\dots$ ---
it oscillates and never converges. The trouble is \emph{period}.
\begin{definition}[Period; aperiodic]
\label{ch24:def-aperiodic}
The \textbf{period} of state $i$ is $d(i) = \gcd\set{n \ge 1 : (P^n)_{ii} > 0}$,
the greatest common divisor of the lengths of all loops returning to $i$. If
$d(i)=1$ the state is \textbf{aperiodic}. In an irreducible chain all states share
the same period; the chain is aperiodic if that common period is $1$. A single
self-loop $P_{ii}>0$ anywhere forces aperiodicity.
\end{definition}
\begin{intuitionbox}
In periodic chains, probability mass is constrained by cyclic timing. In the two-state flip-flop chain,
the time-average fraction of steps spent at state $A$ is $1/2$, but at any single step $n$,
all probability mass resides entirely at one of the two states, alternating with the parity of $n$.
Aperiodicity eliminates this rigid cyclic dependence: when return times can take coprime lengths
(or include self-loops), probability mass disperses across all states simultaneously,
enabling the distribution itself to converge.
\end{intuitionbox}
\subsection{The theorem}
A chain that is both irreducible and aperiodic on a finite state space is called
\textbf{ergodic}. For such chains everything works.
\begin{theorem}[Fundamental theorem of Markov chains / convergence]
\label{ch24:thm-convergence}
Let $P$ be the transition matrix of a finite, irreducible, aperiodic Markov chain.
Then:
\begin{enumerate}
  \item There is a \emph{unique} stationary distribution $\pi$, with $\pi_i > 0$
    for every state $i$.
  \item For \emph{every} starting distribution $\pi^{(0)}$,
    \[
      \pi^{(0)} P^n \longrightarrow \pi \qquad \text{as } n\to\infty .
    \]
    Equivalently, every row of $P^n$ converges to $\pi$: the chain forgets where
    it started.
  \item \textbf{(Ergodic / time-average form.)} With probability one, the
    long-run fraction of time spent in state $i$ equals $\pi_i$:
    \[
      \frac{1}{n}\sum_{t=1}^{n} \ind\{X_t = i\} \asto \pi_i .
    \]
    The stationary mass $\pi_i$ also gives the average wait between returns:
    $\pi_i = 1/\E_i[T_i]$, the reciprocal of the mean return time to state $i$.
\end{enumerate}
\end{theorem}
Parts (2) and (3) are different statements, and the distinction matters. Part (2)
is about the \emph{distribution} at a single far-off time $n$ --- ``where is the
chain likely to be on day 1000?'' Part (3) is about a \emph{single sample path}
--- ``what fraction of its life does this one particular trajectory spend in state
$i$?'' Aperiodicity is needed for (2) (otherwise the distribution oscillates) but
\emph{not} for (3): the flip-flop chain still spends half its time in each state,
even though its distribution never converges. This is why simulations work: a
running average over one long run estimates $\pi$ even for periodic chains.
\begin{whybox}
\textbf{Why does convergence happen, intuitively?} Here is the cleanest argument,
the \emph{coupling} idea. Run two copies of the chain in parallel: copy $X$
started from your distribution of interest, copy $Y$ started from the stationary
$\pi$ (so $Y$ stays distributed as $\pi$ for all time). Let them move
independently until the first moment they happen to land on the same state ---
call it the coupling time $\tau$. After $\tau$, glue them: let $X$ copy $Y$'s moves
exactly. From that instant on, $X$ and $Y$ are identical, so $X$ is also
distributed as $\pi$.
Because the chain is irreducible and aperiodic, the two copies \emph{will} meet
with probability one, and in fact the chance they have \emph{not} met by time $n$
decays geometrically. The total-variation distance between $X$'s distribution and
$\pi$ is bounded by $\P(\tau > n)$, which $\to 0$. So the only obstruction to
convergence is the two copies never managing to meet --- which is exactly what
reducibility (they live in different worlds) or periodicity (they march in
opposite phase and keep missing) would cause. Rule those out and convergence is
forced.
\end{whybox}
\begin{historybox}
Following Andrey Markov's 1906--1913 introduction of dependent sequences,
mathematicians turned to the central question of this chapter: under what
conditions does a chain forget its past and settle into a unique stationary
equilibrium?

Over the 1930s, the general convergence theory was placed on rigorous foundations
by Andrey Kolmogorov and Wolfgang Doeblin. Doeblin, a brilliant French-German
mathematician who tragically died during World War II, invented the revolutionary
coupling technique introduced above: proving convergence without cumbersome matrix
powers simply by racing two copies of a chain until they collide. In parallel, the
spectral foundations were supplied by the Perron--Frobenius theorem of linear
algebra, which guaranteed that irreducible, aperiodic transition matrices have a
simple dominant eigenvalue $\lambda_1 = 1$ with a strictly positive eigenvector $\pi$.

Today, deciding \emph{whether} a finite chain settles is completely understood.
But determining \emph{how fast} it settles in high-dimensional systems remains an
active research frontier: characterizing the sharp cutoff phenomenon (where a chain
hovers unmixed before plunging to equilibrium in just a few steps, as in riffle card
shuffling) and establishing tight spectral gaps and mixing times for complex physical
lattices, spin systems, and randomized algorithms.
\end{historybox}
\begin{figure}[H]
  \centering
  \begin{tikzpicture}
    \begin{axis}[
        width=12cm, height=6cm,
        xlabel={step $t$}, ylabel={empirical fraction of time sunny},
        xmin=0, xmax=400, ymin=0, ymax=1.0,
        title style={font=\small},
        title={Ergodic Theorem: empirical time average along a single sample path},
        grid=both, grid style={gray!18},
        legend pos=north east, legend cell align=left,
      ]
      \addplot[sourceblue,thick] coordinates {
        (1,1)(5,0.8)(10,0.7)(20,0.6)(30,0.633)(45,0.71)(60,0.683)(80,0.65)
        (110,0.672)(150,0.66)(200,0.671)(260,0.663)(320,0.6685)(400,0.6664)};
      \addlegendentry{running fraction in $S$}
      \addplot[black,dashed,thick,domain=0:400] {0.6666667};
      \addlegendentry{$\pi_S = 2/3$}
    \end{axis}
  \end{tikzpicture}
  \caption{The ergodic theorem in action (a single simulated path of the weather
  chain). The fraction of days that are sunny wanders at first, then locks onto
  the stationary value $2/3$. This is part (3) of
  Theorem~\ref{ch24:thm-convergence}: time averages converge along one
  trajectory.}
  \label{ch24:fig-ergodic-timeavg}
\end{figure}
\begin{paradoxbox}
\textbf{Equilibrium without ever being at rest.} Newcomers often expect a chain
``at its stationary distribution'' to slow down or stop. It does the opposite ---
it keeps hopping just as energetically as before. The paradox dissolves once you
separate two ensembles. Imagine a million towns, each running its own weather
chain. If today the \emph{collection} has $2/3$ sunny towns, then tomorrow it
still has $2/3$ sunny towns --- even though roughly $20\%$ of today's sunny towns
turned rainy and exactly enough rainy ones turned sunny to replace them. No single
town is frozen; the \emph{population} is. Confusing the fate of one walker with
the shape of the crowd is the root of nearly every Markov-chain misconception.
\end{paradoxbox}
\section{Reversibility and Detailed Balance}
\label{ch24:reversibility}
Solving $\pi P = \pi$ means solving a linear system --- fine for three states,
miserable for three million. There is a far easier sufficient condition that, when
it holds, hands you $\pi$ almost for free. It is also the single most important
idea in all of MCMC.
The stationarity equation $\pi_j = \sum_i \pi_i P_{ij}$ is a statement about
\emph{total} flow: the net probability current into each state is zero. Detailed
balance demands something stronger and far more rigid --- that the flow balance
between \emph{every individual pair} of states, edge by edge.
\begin{definition}[Detailed balance; reversibility]
\label{ch24:def-detailed-balance}
A distribution $\pi$ satisfies \textbf{detailed balance} for $P$ if
\[
  \pi_i\, P_{ij} = \pi_j\, P_{ji} \qquad \text{for all states } i,j .
\]
A chain that has such a $\pi$ is called \textbf{reversible}.
\end{definition}
The left side is the rate of $i\to j$ traffic in equilibrium; the right side is
the rate of $j\to i$ traffic. Detailed balance says these are equal for each edge:
no net circulation anywhere.
\begin{theorem}[Detailed balance implies stationary]
\label{ch24:thm-db-implies-stationary}
If $\pi$ satisfies detailed balance for $P$, then $\pi$ is stationary: $\pi P = \pi$.
\end{theorem}
\begin{proof}
Sum the detailed-balance equations over $i$ for a fixed $j$:
\[
  \sum_i \pi_i P_{ij} = \sum_i \pi_j P_{ji} = \pi_j \sum_i P_{ji} = \pi_j \cdot 1 = \pi_j ,
\]
using that the $j$-th row of $P$ sums to one. But $\sum_i \pi_i P_{ij} = (\pi P)_j$,
so $(\pi P)_j = \pi_j$ for every $j$, i.e.\ $\pi P = \pi$.
\end{proof}
\begin{happeningbox}
\textbf{What is really happening: time looks the same run backwards.} The name
``reversible'' is literal. Record a long stationary run of the chain and play the
film backwards. If detailed balance holds, the reversed film is statistically
\emph{indistinguishable} from a forward run --- the same transition probabilities
govern it. Check: the probability of seeing $i$ then $j$ forward is
$\pi_i P_{ij}$; the probability of seeing $j$ then $i$ (which is what
``$i$ then $j$'' looks like in reverse) is $\pi_j P_{ji}$. Detailed balance is
exactly the statement that these are equal. So reversibility $=$ time-symmetry of
the equilibrium dynamics. Chains \emph{without} it have a built-in arrow of time:
a preferred direction of circulation, like a whirlpool, that you would notice if
the film were reversed.
\end{happeningbox}

\begin{recoverybox}
\textbf{Pause and think: the arrow of time.} Can a one-way cycle $1 \to 2 \to 3 \to 1$ (where $P_{12}=P_{23}=P_{31}=1$) ever satisfy detailed balance?
\emph{Ten-second thought:} No. Across the edge connecting states $1$ and $2$, forward traffic is $\pi_1 P_{12} = \tfrac13 \cdot 1 = \tfrac13$. But reverse traffic is $\pi_2 P_{21} = \tfrac13 \cdot 0 = 0$. Detailed balance demands that traffic in both directions balances across every individual edge. A one-way street has an unmistakable arrow of time: played in reverse, the film shows walkers moving backward around the cycle.
\end{recoverybox}

\begin{figure}[H]
\centering
\begin{tikzpicture}[
  mystate/.style={circle, draw=#1, thick, fill=#1!12, minimum size=0.9cm, font=\small\bfseries},
  netflux/.style={->, >=Stealth, very thick}
]
  % Left panel: Global Balance (Non-reversible)
  \begin{scope}[xshift=0cm]
    \node[font=\small\bfseries, align=center] at (1.8, 3.4) {Global Balance Only\\ \footnotesize (Net clockwise circulation)};
    \node[mystate=sourceblue] (A) at (1.8, 2.4) {$1$};
    \node[mystate=sourceblue] (B) at (0.2, 0.3) {$2$};
    \node[mystate=sourceblue] (C) at (3.4, 0.3) {$3$};

    \draw[netflux, sourceblue] (A) edge[bend right=22] node[left=2pt, font=\scriptsize] {$0.7$} (B);
    \draw[->, gray!70] (B) edge[bend right=10] node[pos=0.35, right=2pt, font=\tiny] {$0.3$} (A);

    \draw[netflux, sourceblue] (B) edge[bend right=22] node[below=1pt, font=\scriptsize] {$0.7$} (C);
    \draw[->, gray!70] (C) edge[bend right=10] node[pos=0.5, above=1pt, font=\tiny] {$0.3$} (B);

    \draw[netflux, sourceblue] (C) edge[bend right=22] node[right=2pt, font=\scriptsize] {$0.7$} (A);
    \draw[->, gray!70] (A) edge[bend right=10] node[pos=0.65, left=2pt, font=\tiny] {$0.3$} (C);

    \node[font=\scriptsize, align=center, anchor=north] at (1.8, -0.55) {
      $\text{Inflow} = \text{Outflow} = 1/3$
    };
    \node[font=\scriptsize, align=center, anchor=north] at (1.8, -0.95) {
      \textbf{Net edge flux $\ne 0$}
    };
    \node[font=\scriptsize, align=center, anchor=north] at (1.8, -1.45) {
      $\pi_1 P_{12} = \tfrac{7}{30} \ne \pi_2 P_{21} = \tfrac{3}{30}$
    };
  \end{scope}

  % Divider
  \draw[dashed, gray!40, thick] (4.9, -1.8) -- (4.9, 3.6);

  % Right panel: Detailed Balance (Reversible)
  \begin{scope}[xshift=6.4cm]
    \node[font=\small\bfseries, align=center] at (1.8, 3.4) {Detailed Balance\\ \footnotesize (Zero net edge current)};
    \node[mystate=targetgreen] (X) at (1.8, 2.4) {$1$};
    \node[mystate=targetgreen] (Y) at (0.2, 0.3) {$2$};
    \node[mystate=targetgreen] (Z) at (3.4, 0.3) {$3$};

    \draw[<->, >=Stealth, very thick, targetgreen] (X) -- node[midway, left=2pt, font=\scriptsize, align=right] {$\pi_1 P_{12} = \pi_2 P_{21}$\\ $\tfrac{1}{6} = \tfrac{1}{6}$} (Y);
    \draw[<->, >=Stealth, very thick, targetgreen] (Y) -- node[midway, below=3pt, font=\scriptsize] {$\pi_2 P_{23} = \pi_3 P_{32} = \tfrac{1}{6}$} (Z);
    \draw[<->, >=Stealth, very thick, targetgreen] (X) -- node[midway, right=2pt, font=\scriptsize, align=left] {$\pi_1 P_{13} = \pi_3 P_{31}$\\ $\tfrac{1}{6} = \tfrac{1}{6}$} (Z);

    \node[font=\scriptsize, align=center, anchor=north] at (1.8, -0.6) {
      \textbf{Zero net flux on every pipe}\\
      Forward and backward movies\\
      statistically identical
    };
  \end{scope}

\end{tikzpicture}
\caption{Global balance versus detailed balance. Left: under global balance ($\pi P = \pi$), total probability entering each state equals total probability leaving, yet a steady circulating current swirls around the loop (here clockwise with net flow $\tfrac{4}{30}$). Right: under detailed balance ($\pi_i P_{ij} = \pi_j P_{ji}$), the traffic across every individual edge balances in both directions, eliminating all net circulation and establishing time-reversibility.}
\label{ch24:fig-balance-comparison}
\end{figure}

\begin{example}[Random walk on a graph]
\label{ch24:ex-graph-walk}
Take any connected undirected graph. A walker sits at a vertex and at each step
moves to a uniformly random neighbor. If vertex $i$ has degree $d_i$ (number of
edges), then $P_{ij} = 1/d_i$ when $i,j$ are adjacent, else $0$. Guess
$\pi_i \propto d_i$, i.e.\ $\pi_i = d_i / (2E)$ where $E$ is the number of edges
(so $\sum_i d_i = 2E$). Check detailed balance on an edge $\{i,j\}$:
\[
  \pi_i P_{ij} = \frac{d_i}{2E}\cdot\frac{1}{d_i} = \frac{1}{2E}
  = \frac{d_j}{2E}\cdot\frac{1}{d_j} = \pi_j P_{ji}. \checkmark
\]
It balances for every edge, so $\pi_i = d_i/(2E)$ is stationary. We \emph{guessed}
the answer and verified it in one line --- no linear system. The lesson: a random
walker lingers in proportion to how connected a place is. High-degree hubs are
visited more. This single observation is the seed of PageRank.
\end{example}

\begin{aestheticbox}
Detailed balance is the probabilist's version of a conservation law in physics ---
indeed it \emph{is} one. In statistical mechanics the equilibrium (Boltzmann)
distribution $\pi_i \propto e^{-E_i/kT}$ satisfies detailed balance with respect
to the microscopic dynamics, and that is precisely why thermal systems settle. The
same equation that tells you the weather will be sunny $2/3$ of the time tells a
physicist how molecules distribute among energy levels. When you reach
Section~\ref{ch24:mcmc} you will see we run the logic constructively: given a
target distribution $\pi$, we \emph{design} a transition matrix $P$ that satisfies detailed balance,
ensuring the chain converges to the intended target distribution.
\end{aestheticbox}

\section{Mixing Time: How Long Until ``Long Run''?}
\label{ch24:mixing}

Theorem~\ref{ch24:thm-convergence} promises convergence but is silent about
\emph{speed}. In Figure~\ref{ch24:fig-weather-converge} we were essentially done
in five days. A poorly connected chain on a million states might need millions of
steps. ``How long?'' is the question of \textbf{mixing time}, and for
applications it is everything: a sampler that has not yet mixed gives you biased
garbage.

To measure ``how far from $\pi$,'' we need a distance between distributions.

\begin{definition}[Total variation distance]
\label{ch24:def-tv}
For two distributions $\mu,\nu$ on the same states,
\[
  \norm{\mu - \nu}_{\mathrm{TV}} = \tfrac12 \sum_i \abs{\mu_i - \nu_i}
  = \max_{A} \abs{\mu(A) - \nu(A)} .
\]
It is $0$ when the distributions coincide and $1$ when their supports are disjoint;
the second form says it is the largest gap in probability they assign to any event.
\end{definition}

\begin{definition}[Mixing time]
\label{ch24:def-mixing}
Fix a tolerance $\eps$ (commonly $\eps = 1/4$). The \textbf{mixing time} is
\[
  t_{\mathrm{mix}}(\eps) = \min\set[\Big]{ n : \max_{x}\ \norm{\,\delta_x P^n - \pi\,}_{\mathrm{TV}} \le \eps },
\]
the number of steps after which, \emph{no matter the starting state} $x$, the
chain's distribution is within $\eps$ of stationary in total variation.
\end{definition}

\subsection{The spectral gap controls the rate}

Here the eigenvalues come back to do real work. Order the eigenvalues of $P$ by
modulus: $1 = \lambda_1 > \abs{\lambda_2} \ge \cdots \ge \abs{\lambda_m}$. For an
ergodic chain the top one is $1$ (the stationary direction) and all others are
strictly smaller. Decompose any starting vector along the eigen-directions; under
repeated multiplication by $P$, the stationary component is preserved while every
other component is multiplied by its $\lambda_k$ each step and \emph{decays
geometrically}. The slowest-dying piece is governed by the second eigenvalue.

\begin{theorem}[Geometric convergence rate]
\label{ch24:thm-spectral}
For an ergodic chain (reversible, say, so the $\lambda_k$ are real), there is a
constant $C$ with
\[
  \norm{\,\delta_x P^n - \pi\,}_{\mathrm{TV}} \le C\,\abs{\lambda_2}^{\,n}.
\]
Writing the \textbf{spectral gap} as $\gamma = 1 - \abs{\lambda_2}$, the mixing
time scales like
\[
  t_{\mathrm{mix}} \approx \frac{1}{\gamma}\,\log\frac{1}{\eps}.
\]
A big gap (small $\abs{\lambda_2}$) means fast mixing; a tiny gap means a sluggish
chain.
\end{theorem}

For the weather chain the eigenvalues of $P$ are $\lambda_1 = 1$ and
$\lambda_2 = 0.8 + 0.6 - 1 = 0.4$ (the eigenvalues of a $2\times2$ stochastic
matrix are $1$ and $\mathrm{tr}\,P - 1$). The gap is $0.6$ --- large --- so error
falls by a factor of $0.4$ per day, matching the rapid collapse in
Figure~\ref{ch24:fig-weather-converge}. Indeed the gap between the two curves
there shrinks by exactly $0.4$ each step: that is $\abs{\lambda_2}^n$ in the flesh.

\begin{figure}[H]
  \centering
  \begin{tikzpicture}
    \begin{axis}[
        width=12cm, height=6cm,
        xlabel={step $n$}, ylabel={TV distance to $\pi$ (log scale)},
        ymode=log, xmin=0, xmax=12, ymin=1e-5, ymax=1,
        title style={font=\small},
        title={Total-variation mixing rate governed by the second eigenvalue $\abs{\lambda_2}$},
        grid=both, grid style={draw=gray!50, opacity=0.2},
        legend style={at={(0.05, 0.05)}, anchor=south west, column sep=6pt, font=\small}, legend cell align=left,
      ]
      \addplot[sourceblue,very thick,mark=*,mark size=1.3pt,domain=0:12,samples=13] {0.5*0.4^x};
      \addlegendentry{fast: $\abs{\lambda_2}=0.4$}
      \addplot[densityred,very thick,mark=square*,mark size=1.3pt,domain=0:12,samples=13] {0.5*0.9^x};
      \addlegendentry{slow: $\abs{\lambda_2}=0.9$}
    \end{axis}
  \end{tikzpicture}
  \caption{On a log scale, total-variation distance to $\pi$ falls along a
  straight line whose slope is $\log\abs{\lambda_2}$. The fast chain (steep line)
  mixes in a few steps; the slow chain (shallow line) crawls. The spectral gap governs the slope.}
  \label{ch24:fig-mixing-rate}
\end{figure}

\begin{intuitionbox}
What makes a chain slow to mix is a \textbf{bottleneck}: a narrow bridge between
two roomy regions of state space (Figure~\ref{ch24:fig-bottleneck}). The walker
spends ages on one side, only rarely squeezing through. Until it crosses, it
cannot ``know'' about the far side, so its distribution stays lopsided.
Mathematically, a bottleneck of small ``conductance'' (restricted outflow relative
to room size) forces $\abs{\lambda_2}$ close to $1$ and the gap close to $0$. This
is a primary challenge in MCMC: a multimodal target with widely separated modes
creates a severe bottleneck in state space, where a naive sampler can remain trapped in
one mode for many iterations and fail to visit the others.
\end{intuitionbox}

\begin{figure}[H]
  \centering
  \begin{tikzpicture}[>=Stealth,
    roomnode/.style={circle, draw=#1, very thick, fill=#1!15, minimum size=0.85cm, font=\small\bfseries},
    bridgenode/.style={circle, draw=#1, very thick, fill=#1!25, minimum size=0.9cm, font=\small\bfseries}
  ]
    \draw[sourceblue!40, dashed, thick, fill=sourceblue!5, rounded corners=10pt]
      (-0.3, -1.7) rectangle (3.3, 1.7);
    \node[font=\small\bfseries, sourceblue!85!black] at (1.5, 1.35) {Room 1 (Fast local mixing)};

    \draw[targetgreen!40, dashed, thick, fill=targetgreen!5, rounded corners=10pt]
      (6.7, -1.7) rectangle (10.3, 1.7);
    \node[font=\small\bfseries, targetgreen!85!black] at (8.5, 1.35) {Room 2 (Fast local mixing)};

    \node[roomnode=sourceblue] (A1) at (0.8, 0.55) {$1$};
    \node[roomnode=sourceblue] (A2) at (0.8, -0.75) {$2$};
    \node[bridgenode=sourceblue] (A3) at (2.4, 0.0) {$3$};

    \draw[<->, thick, sourceblue] (A1) -- node[left=1pt, font=\scriptsize] {$0.4$} (A2);
    \draw[<->, thick, sourceblue] (A1) -- node[above=1pt, font=\scriptsize] {$0.5$} (A3);
    \draw[<->, thick, sourceblue] (A2) -- node[below=1pt, font=\scriptsize] {$0.5$} (A3);

    \node[bridgenode=targetgreen] (B1) at (7.6, 0.0) {$4$};
    \node[roomnode=targetgreen] (B2) at (9.2, 0.55) {$5$};
    \node[roomnode=targetgreen] (B3) at (9.2, -0.75) {$6$};

    \draw[<->, thick, targetgreen] (B2) -- node[right=1pt, font=\scriptsize] {$0.4$} (B3);
    \draw[<->, thick, targetgreen] (B1) -- node[above=1pt, font=\scriptsize] {$0.5$} (B2);
    \draw[<->, thick, targetgreen] (B1) -- node[below=1pt, font=\scriptsize] {$0.5$} (B3);

    \draw[<->, >=Stealth, very thick, densityred] (A3) --
      node[midway, above=3pt, font=\small\bfseries] {rare bridge: $\eps \ll 1$}
      node[midway, below=3pt, font=\scriptsize, align=center] {bottleneck flow\\ $\pi_3 P_{34} = \pi_4 P_{43} \approx 0$} (B1);
  \end{tikzpicture}
  \caption{The geometry of slow mixing: a bottleneck in state space. Within each cluster of states (Room 1 and Room 2), transitions are frequent and the walker mixes locally in just a few steps. But crossing the bridge between states $3$ and $4$ has a tiny probability $\eps \ll 1$. The chain spends long epochs trapped on one side before rare chance lets it cross, forcing the second eigenvalue $\abs{\lambda_2}$ close to $1$ (a tiny spectral gap $\gamma \approx 2\eps$) and driving the mixing time up to order $1/\eps$.}
  \label{ch24:fig-bottleneck}
\end{figure}

\begin{paradoxbox}
\textbf{The cutoff phenomenon.} One might expect mixing to be gradual, with total-variation
distance decaying steadily at each step. For many natural chains, however, the
decay exhibits a sharp \emph{cutoff}. Persi Diaconis and collaborators showed that
riffle-shuffling a $52$-card deck behaves this way: after $5$ or $6$ shuffles the
total-variation distance to uniform is still close to $1$, around $7$ shuffles it plunges,
and by $8$ or $9$ shuffles the deck is well mixed. More generally, in chains with cutoff,
the distance to stationarity stays near $1$ up to a sharp threshold and then drops to near
$0$ over a narrow window.
\end{paradoxbox}

\section{Headline Application I: PageRank}
\label{ch24:pagerank}

\begin{applicationbox}
\textbf{Google's original ranking algorithm is a stationary distribution.} Model
the web as a giant directed graph: pages are states, links are edges. Imagine a
\emph{random surfer} who, at each page, clicks a uniformly random outgoing link.
This is a Markov chain with $P_{ij} = 1/(\text{outdegree of }i)$ for each link
$i\to j$. The PageRank of a page is defined as its stationary probability $\pi_i$:
the long-run fraction of time the surfer spends there. Pages that many
(important) pages link to get visited often, so they score high --- importance
defined self-referentially and resolved by the eigenvector: PageRank is
precisely \emph{the left eigenvector of the web's transition matrix for eigenvalue $1$}
(Section~\ref{ch24:stationary}). Brin and Page's 1998 insight was that this single
number, computable for the whole web, ranks pages very well.

\textbf{Two repairs are needed} to guarantee a unique limit
(Theorem~\ref{ch24:thm-convergence}). (1) \emph{Dangling pages} with no outgoing
links are dead ends --- the surfer gets stuck, breaking irreducibility. (2) The
web is not strongly connected; it has many traps. Both are fixed by the
\textbf{teleportation} trick. With probability $d \approx 0.85$ the surfer
follows a random link as before; with probability $1-d$ she jumps to a page chosen
uniformly at random from all $N$ pages. The modified matrix is
\[
  G = d\,P + (1-d)\,\tfrac{1}{N}\,\mathbf{J},
\]
where $\mathbf{J}$ is the all-ones matrix (the ``Google matrix''). Because $G$ has
every entry strictly positive, the chain is automatically irreducible and
aperiodic, so a unique positive $\pi$ exists and the surfer converges to it from
any start. PageRank is just $\pi G = \pi$.
\end{applicationbox}

\begin{happeningbox}
\textbf{How is $\pi$ actually computed for billions of pages?} Not by solving a
linear system (infeasible at web scale) and not by finding eigenvectors with
standard linear-algebra routines. It is computed by \emph{the power iteration of
Section~\ref{ch24:longrun}}: pick any starting vector, repeatedly multiply by $G$,
and watch it converge. This is exactly $\pi^{(n)} = \pi^{(0)} G^n$. The damping
factor $d=0.85$ sets the second eigenvalue to about $0.85$, so the spectral gap is
about $0.15$ and the error drops by $0.85$ per iteration --- roughly $50$--$100$
iterations reach acceptable accuracy. The whole engine that organized the early
web is the convergence theorem plus a matrix-vector product run a few dozen times.
This is why the eigenvalue framing in this chapter is not decoration: it
\emph{is} the algorithm.
\end{happeningbox}

\begin{example}[A tiny PageRank]
\label{ch24:ex-pagerank}
Reuse the three-page chain of Example~\ref{ch24:ex-three-state} (pretend it is
already strongly connected, so we can skip teleportation). We found
$\pi = (5/19,\,8/19,\,6/19) \approx (0.263,\,0.421,\,0.316)$. So page $B$ ranks
first, $C$ second, $A$ third. Notice the ranking is \emph{not} simply ``count the
inbound links'': it weights each inbound link by the importance of its source and
splits a page's vote across its outbound links. That recursive weighting is the
whole point, and the stationary distribution resolves the recursion in one shot.
\end{example}

\section{Headline Application II: Markov Chain Monte Carlo}
\label{ch24:mcmc}

\begin{applicationbox}
\textbf{Sampling from impossible distributions by building a chain on purpose.}
In Bayesian statistics, machine learning, and physics you constantly need to draw
samples from a distribution $\pi$ known only up to a constant:
$\pi(x) = \tilde\pi(x)/Z$, where $\tilde\pi$ is computable but the normalizer
$Z = \sum_x \tilde\pi(x)$ is a sum over astronomically many states (every
parameter setting, every spin configuration). You cannot compute $Z$, so you
cannot use it directly --- yet you want samples from $\pi$.

\textbf{MCMC turns the whole chapter inside out.} Instead of being handed a chain
and asked for its $\pi$, you are handed a target $\pi$ and you \emph{construct a
chain whose stationary distribution is that $\pi$}. Run the chain long enough
(past its mixing time) and its state is an (approximate) sample from $\pi$.
Convergence (Theorem~\ref{ch24:thm-convergence}) guarantees this works; detailed
balance (Section~\ref{ch24:reversibility}) is how you build it; mixing time
(Section~\ref{ch24:mixing}) tells you how long to wait.

\textbf{Division of labor with Chapter 35.} In this section, we study MCMC in its
purest, finite-state discrete form to see detailed balance and the cancellation of $Z$
with absolute mathematical clarity. In Chapter 35 (Monte Carlo Methods), we expand this
seed into the full continuous machinery: continuous proposal densities $q(y\given x)$,
Gibbs sampling, and industrial diagnostics for high-dimensional models.
\end{applicationbox}

The standard method is the \textbf{Metropolis--Hastings} algorithm, which uses detailed
balance constructively. We \emph{want} $\pi_i P_{ij} = \pi_j P_{ji}$; we
\emph{design} $P$ to make it so. The recipe: from state $i$, \emph{propose} a move
to a candidate $j$ using any convenient proposal $Q_{ij}$ (say, perturb the
current parameters a little). Then \emph{accept} the move with probability
\[
  \alpha_{ij} = \min\set[\Big]{1,\ \frac{\tilde\pi_j\, Q_{ji}}{\tilde\pi_i\, Q_{ij}}};
\]
otherwise stay at $i$. The realized transition probability is
$P_{ij} = Q_{ij}\,\alpha_{ij}$ for $j\neq i$.

\begin{whybox}
\textbf{Why the acceptance rule produces the right $\pi$.} Check detailed balance
directly. Suppose, without loss of generality, that the ratio
$r = \tilde\pi_j Q_{ji} / (\tilde\pi_i Q_{ij}) \le 1$, so $\alpha_{ij} = r$ and
$\alpha_{ji} = 1$. Then
\[
  \pi_i P_{ij} = \pi_i\, Q_{ij}\,\alpha_{ij}
  = \frac{\tilde\pi_i}{Z}\, Q_{ij}\cdot \frac{\tilde\pi_j Q_{ji}}{\tilde\pi_i Q_{ij}}
  = \frac{\tilde\pi_j Q_{ji}}{Z},
\]
while
\[
  \pi_j P_{ji} = \pi_j\, Q_{ji}\,\alpha_{ji}
  = \frac{\tilde\pi_j}{Z}\, Q_{ji}\cdot 1 = \frac{\tilde\pi_j Q_{ji}}{Z}.
\]
They are equal, so detailed balance holds and $\pi$ is stationary by
Theorem~\ref{ch24:thm-db-implies-stationary}. \emph{The unknown normalizer $Z$
cancels} --- it appears in both sides and divides out, and inside $\alpha$ only the
ratio $\tilde\pi_j/\tilde\pi_i$ ever appears. This exact cancellation makes MCMC practical:
it generates samples from $\pi$ without requiring evaluation of the normalizing constant $Z$.
\end{whybox}

\begin{labbox}
Convince yourself the cancellation is real, on a baby example. Let the target be
$\tilde\pi = (1, 3, 2)$ on three states (so the true $\pi = (1/6, 1/2, 1/3)$, but
pretend you do not know the denominator $6$). Use a symmetric proposal: from any
state, propose each of the other two with probability $1/2$, so $Q_{ij}=Q_{ji}$
and the acceptance ratio is just $\tilde\pi_j/\tilde\pi_i$. Work out the
acceptance probability for the proposal $1 \to 2$ (it is
$\min\{1, 3/1\} = 1$, always accept --- you are climbing toward higher
probability) and for $2 \to 1$ (it is $\min\{1, 1/3\} = 1/3$ --- you usually
refuse to descend). Simulate a few hundred steps by hand or by code and tally the
visits; the fractions should approach $(1/6, 1/2, 1/3)$ (Figure~\ref{ch24:fig-mcmc-hist}).
You will have sampled from $\pi$ using only the \emph{ratios} of $\tilde\pi$.
\end{labbox}

\begin{figure}[H]
  \centering
  \begin{tikzpicture}
    \begin{axis}[
        width=11cm, height=6cm,
        ybar, bar width=14pt,
        xlabel={state}, ylabel={probability},
        symbolic x coords={1,2,3},
        xtick=data, ymin=0, ymax=0.6,
        enlarge x limits=0.3,
        title style={font=\small},
        title={Metropolis sampler on unnormalized target $\tilde\pi=(1,3,2)$},
        legend pos=north west, legend cell align=left,
        grid=both, grid style={gray!18},
      ]
      \addplot[draw=sourceblue,fill=sourceblue!35] coordinates {(1,0.1667)(2,0.5)(3,0.3333)};
      \addlegendentry{target $\pi$}
      \addplot[draw=densityred,fill=densityred!35] coordinates {(1,0.169)(2,0.500)(3,0.331)};
      \addlegendentry{Metropolis (20000 steps)}
    \end{axis}
  \end{tikzpicture}
  \caption{The Metropolis sampler of the lab box in action on
  $\tilde\pi=(1,3,2)$. Tallying the visited states over $20{,}000$ steps
  reproduces the target $\pi=(1/6,1/2,1/3)$ to within sampling noise --- all
  while never computing the normalizer $Z=6$. The chain samples the right shape
  from the ratios alone.}
  \label{ch24:fig-mcmc-hist}
\end{figure}

\begin{connectionbox}
MCMC is the conceptual bridge from this chapter to Bayesian inference (Chapter 31)
and computational Monte Carlo (Chapter 35). In Bayesian inference, the target is the
posterior $\pi(\theta) \propto \text{likelihood}(\theta)\times \text{prior}(\theta)$,
whose normalizing constant --- the marginal likelihood or ``evidence'' --- is an
intractable high-dimensional integral.

Here you have seen the core discrete engine: detailed balance guarantees that $\pi$
is stationary, while the unknown normalizer $Z$ cancels out in the acceptance ratio
$\alpha$. In Chapter 35, we will build out the full computational system: transitioning
from discrete state spaces to continuous distributions, replacing proposal matrices
with continuous densities, introducing Gibbs sampling, and using diagnostics
(trace plots, autocorrelation times, and Gelman--Rubin $\hat{R}$) to test whether
the mixing time of Section~\ref{ch24:mixing} has genuinely passed.
\end{connectionbox}

\begin{trapbox}
\textbf{Burn-in and the perils of not-yet-mixed chains.} A freshly started MCMC
chain reflects its (arbitrary) starting point, not $\pi$. Samples from the early
phase are biased, so practitioners discard an initial \emph{burn-in} segment ---
an attempt to wait out the mixing time. But there is no reliable way to be
\emph{certain} a chain has mixed from the inside: a chain trapped in one mode of a
multimodal target looks perfectly converged (flat trace, stable averages) while
completely missing the other modes. Convergence diagnostics reduce the risk but
cannot eliminate it. The spectral-gap intuition is the right mental safeguard: if
you suspect your target has well-separated regions, suspect a bottleneck, and
suspect your chain.
\end{trapbox}

\section{Worked Examples}
\label{ch24:worked}

\begin{example}[Stationary distribution of a $3\times 3$ chain]
\label{ch24:ex-3x3}
Let
\[
  P = \begin{pmatrix}
    0.5 & 0.5 & 0 \\
    0.25 & 0.5 & 0.25 \\
    0 & 0.5 & 0.5
  \end{pmatrix}.
\]
This is a random walk on a path $1 - 2 - 3$ with self-loops --- irreducible
(you can get anywhere) and aperiodic (self-loops). Solve $\pi P = \pi$ with
$\pi = (a,b,c)$:
\[
  a = 0.5a + 0.25b,\qquad c = 0.25b + 0.5c .
\]
The first gives $0.5a = 0.25b$, so $b = 2a$; the third gives $0.5c = 0.25b$, so
$b = 2c$, hence $a = c$. Normalizing $a + 2a + a = 4a = 1$ gives
$\pi = (1/4,\,1/2,\,1/4)$. \emph{Sanity check via detailed balance:} this is a
path graph; treating self-loops as not affecting the edge balance, the chain is
reversible, and indeed $\pi_1 P_{12} = \tfrac14\cdot\tfrac12 = \tfrac18 =
\tfrac12\cdot\tfrac14 = \pi_2 P_{21}$. \checkmark
\end{example}

\begin{example}[Mean return time]
\label{ch24:ex-return-time}
In the weather chain, how many days on average between rainy days? By part (3) of
Theorem~\ref{ch24:thm-convergence}, the mean return time to a state is
$1/\pi_i$. Since $\pi_R = 1/3$, the expected gap between consecutive rainy days
(counting from one rainy day to the next return) is $1/\pi_R = 3$ days. Likewise a
sunny day recurs on average every $1/\pi_S = 1.5$ days. The rarer state has the
longer return time --- exactly as it should.
\end{example}

\begin{example}[Build a sampler with Metropolis]
\label{ch24:ex-metropolis}
You want samples from $\pi \propto (4, 2, 1)$ on states $\{1,2,3\}$ (true
$\pi = (4/7, 2/7, 1/7)$, pretend $7$ unknown). Use the symmetric proposal ``flip a
coin to propose a neighbor'' on the cycle $1\!-\!2\!-\!3\!-\!1$. Acceptance for
$1\to 2$: $\min\{1, 2/4\} = 1/2$. For $2\to 1$: $\min\{1, 4/2\}=1$. For $2\to 3$:
$\min\{1,1/2\}=1/2$. For $3\to 2$: $\min\{1, 2\}=1$. Detailed balance check on edge
$\{1,2\}$ with $Q_{12}=Q_{21}=1/2$:
\[
  \pi_1 P_{12} = \tfrac47\cdot\tfrac12\cdot\tfrac12 = \tfrac{1}{7},\qquad
  \pi_2 P_{21} = \tfrac27\cdot\tfrac12\cdot 1 = \tfrac{1}{7}. \checkmark
\]
Balanced, so the chain's stationary distribution is the target. The asymmetry of
acceptance --- always climb, sometimes descend --- is what concentrates the
samples on the high-probability state $1$.
\end{example}

\section{Looking Ahead: From Discrete Clicks to Continuous Time}
\label{ch24:looking-ahead}

Throughout Chapters 23 and 24, time has marched to the beat of an external
metronome: step $0, 1, 2, \dots$ At each discrete tick of the clock, the system
observes where it currently sits and chooses where to hop next. That discrete
cadence let us capture the entire future evolution of the chain in the powers
$P^n$ of a single transition matrix.

Yet nature rarely waits for clock ticks. Radioactive atoms do not consult a
stopwatch before decaying; photons strike a telescope detector whenever they
arrive; server requests and incoming calls land along a continuous, unbroken
timeline. What happens when transitions can occur at \emph{any} real instant
$t \ge 0$, while still preserving the memoryless property that made Markov chains
so tractable?

In Chapter 25, we take that step into continuous time with the \textbf{Poisson
process} --- the universal baseline model for memoryless random arrivals
occurring at a constant average rate. Later, in Chapter 27, we will merge
continuous time with the multi-state graphs of this chapter to construct
continuous-time Markov chains and queueing networks. But the bridge begins with
the arrival stream itself.

\section{Exercises}
\label{ch24:exercises}

\begin{exercise}[Two-state mixing rate]
\label{ch24:ex-two-state-mixing}
For the weather matrix $P=\left(\begin{smallmatrix}0.8 & 0.2\\ 0.4 & 0.6\end{smallmatrix}\right)$:
\begin{enumerate}[label=(\alph*)]
  \item \emph{(Mental check.)} Before calculating, explain intuitively why $\pi_S$
    must exceed $1/2$ by comparing the probability of leaving the sun ($P_{SR}=0.2$)
    with the probability of leaving the rain ($P_{RS}=0.4$).
  \item Verify that the eigenvalues of $P$ are $1$ and $0.4$. Starting from $(1,0)$,
    find a closed form for $\pi^{(n)}_S$ and confirm it converges to $2/3$ with the
    error shrinking by a factor of $0.4$ each step. (Hint: write $\pi^{(0)}$ as a
    linear combination of the left eigenvectors.)
\end{enumerate}
\end{exercise}

\begin{exercise}[Find the stationary distribution]
\label{ch24:ex-find-stationary}
Compute the stationary distribution of
\[
  P = \begin{pmatrix} 0.2 & 0.6 & 0.2 \\ 0.3 & 0.4 & 0.3 \\ 0.5 & 0.1 & 0.4 \end{pmatrix}.
\]
Confirm your answer satisfies $\pi P = \pi$ and sums to $1$.
\end{exercise}

\begin{exercise}[Reducible and periodic]
\label{ch24:ex-reducible-periodic}
\begin{enumerate}[label=(\alph*)]
  \item Give an example of an irreducible chain with no self-loops that is periodic,
    and state its period.
  \item Give an example of a chain that has \emph{two} different stationary distributions,
    and explain which hypothesis of Theorem~\ref{ch24:thm-convergence} it violates.
\end{enumerate}
\end{exercise}

\begin{exercise}[Random walk on a graph]
\label{ch24:ex-graph-walk-exercise}
Consider the random walk on the square graph (a $4$-cycle $1\!-\!2\!-\!3\!-\!4\!-\!1$).
Use the degree formula from Example~\ref{ch24:ex-graph-walk} to write down $\pi$
without solving any linear system, then verify $\pi P = \pi$. Is this chain
aperiodic? What goes wrong, and how would a single self-loop fix it?
\end{exercise}

\begin{exercise}[Detailed balance fails]
\label{ch24:ex-detailed-balance-fails}
Show that the cyclic chain
$P=\left(\begin{smallmatrix}0&1&0\\0&0&1\\1&0&0\end{smallmatrix}\right)$ has
stationary distribution $(1/3,1/3,1/3)$ but does \emph{not} satisfy detailed
balance. Interpret the failure in terms of the ``arrow of time'': what does the
chain look like played backwards?
\end{exercise}

\begin{exercise}[Metropolis by hand]
\label{ch24:ex-metropolis-by-hand}
You want to sample from $\tilde\pi = (1, 4, 9, 16)$ on states $\{1,2,3,4\}$ using
the symmetric ``propose a uniform random other state'' proposal.
\begin{enumerate}[label=(\alph*)]
  \item \emph{(Mental check.)} Without computing numbers, which state must the
    chain visit most often, and why does detailed balance force this even if the
    proposal is completely uniform?
  \item Write down the acceptance probability for every ordered pair $(i,j)$, $i\neq j$.
  \item Verify detailed balance on the pair $\{2,4\}$.
\end{enumerate}
\end{exercise}

\begin{exercise}[Spectral gap and bottlenecks]
\label{ch24:ex-bottlenecks}
Consider the bottleneck graph illustrated in Figure~\ref{ch24:fig-bottleneck},
with two dense rooms of $k$ states joined by a single low-probability bridge edge
with transition probability $\eps \ll 1$. Argue qualitatively why the second
eigenvalue $\abs{\lambda_2}$ is close to $1$ and the mixing time is of order $O(1/\eps)$.
Connect this to why a multimodal target distribution is notoriously hard for MCMC samplers.
\end{exercise}

\begin{exercise}[PageRank with teleportation]
\label{ch24:ex-pagerank-teleport}
Take the two-page web $1 \to 2$, $2 \to 1$ (each page links only to the other).
\begin{enumerate}[label=(\alph*)]
  \item Write the raw transition matrix $P$ and its stationary distribution.
  \item Form the Google matrix $G = 0.85\,P + 0.15\cdot\tfrac12 \mathbf{J}$ and find
    its stationary distribution.
  \item Now add a dangling page $3$ with no outlinks; explain why $P$ alone has no
    well-defined surfer chain and how teleportation repairs it.
\end{enumerate}
\end{exercise}

\begin{exercise}[Proof: detailed balance gives stationarity]
\label{ch24:ex-proof-db}
Prove Theorem~\ref{ch24:thm-db-implies-stationary} from first principles directly from
definitions: assuming $\pi_i P_{ij} = \pi_j P_{ji}$ for all $i,j$, show $\pi P = \pi$.
Where exactly does the proof use that the rows of $P$ sum to one?
\end{exercise}

\begin{exercise}[Time average vs.\ distribution]
\label{ch24:ex-time-avg}
For the flip-flop chain $P=\left(\begin{smallmatrix}0&1\\1&0\end{smallmatrix}\right)$
started at state $1$:
\begin{enumerate}[label=(\alph*)]
  \item Does the distribution vector $\pi^{(n)}$ converge?
  \item Does the time-average fraction of visits to state $1$ converge, and to what?
  \item Reconcile your answers with parts (2) and (3) of Theorem~\ref{ch24:thm-convergence},
    identifying which hypothesis fails and which conclusion survives anyway.
\end{enumerate}
\end{exercise}
